The first two sections established that structural aliasing operates at the level of raw signal samples and token vectors.
However, in practical deployment settings, time-series models are not evaluated on abstract sinusoids: they are tasked with recognising \emph{domain concepts} — states, events, anomalies, and transitions that carry operational meaning.
A model with frequency-dependent blind spots may therefore fail not merely in terms of spectral reconstruction accuracy, but in terms of its ability to correctly identify the concepts that downstream decisions depend upon.

This section makes that connection explicit by introducing a formal semantic layer.
We adopt an ontology-based representation~\cite{voitaInformationTheoreticProbingMinimum2020} that articulates the relevant concepts, their properties, and the relations among them in the chosen application domain.
The ontology is constructed in three stages: (i) identification of the core domain concepts (e.g.\ operating regime, fault condition, periodic load cycle) and their characteristic frequency signatures; (ii) formalisation of the relationships between these concepts and the signal-level observability conditions derived in Sections~\ref{sec:one} and~\ref{sec:two}; and (iii) analysis of which concept classes are structurally at risk given the patch parameters of Chronos-Bolt.

\subsection{Ontological Representation}

\noindent\textbf{[To be completed.]}
The ontology will be formally specified here, with class definitions, property axioms, and a mapping from concept identifiers to their expected frequency ranges and temporal granularities.

\subsection{Impact of Structural Aliasing on Semantic Observability}

\noindent\textbf{[To be completed.]}
This subsection will trace how aliasing at blind-spot frequencies propagates upward to distort the recognition of specific ontology classes, providing concrete examples from the chosen domain.
